\section{CONSIDERAÇÕES FINAIS}

Este trabalho teve como objetivo desenvolver uma aplicação móvel colaborativa para apoiar tutores e a comunidade na divulgação e na busca por animais perdidos e encontrados. A versão funcional do WeFIND reúne cadastro de publicações, consulta em mapa, registro de avistamentos, filtros de busca, mensagens privadas e geração de cartazes digitais. A implementação atende ao objetivo de construir e avaliar um protótipo funcional, mas não permite concluir que o aplicativo aumente a quantidade de reencontros ou reduza o tempo necessário para localizar um animal.

Na avaliação de usabilidade, 12 tutores voluntários realizaram um roteiro de cinco tarefas e responderam à Escala de Usabilidade do Sistema (SUS). Todos declararam ter concluído as tarefas com facilidade, e a média SUS foi de 93,33 pontos. Nas respostas abertas, participantes destacaram o mapa, a localização e os recursos de publicação, mas também relataram dúvidas sobre como abrir os detalhes de uma publicação, distinguir as opções de registro e localizar o histórico. Portanto, a avaliação aponta uma percepção favorável entre as pessoas consultadas, ao mesmo tempo que indica ajustes concretos para a interface. O resultado deve ser interpretado como exploratório: a amostra foi pequena e não probabilística, e a pergunta sobre tarefas registrou apenas uma percepção global, sem medir o sucesso de cada atividade separadamente.

A redação em português dos itens SUS foi adotada conforme apresentada por Teixeira (2015). Como essa fonte é uma publicação de divulgação, e não um estudo de validação psicométrica da tradução, o escore é apresentado como resultado do instrumento aplicado neste estudo, sem atribuir à versão utilizada uma validação que não foi estabelecida.

Também foram realizadas uma verificação funcional manual e uma análise estática de segurança. A análise automatizada do APK pelo MobSF apontou itens para revisão, mas seus alertas não foram confirmados individualmente como vulnerabilidades. Como não foram realizados testes dinâmicos, teste de invasão ou avaliação independente da infraestrutura hospedada, essas atividades não permitem afirmar que o sistema esteja livre de falhas. Da mesma forma, os relatos de navegação fluida não substituem medições objetivas de desempenho.

\subsection{Dificuldades Encontradas}

Uma decisão difícil foi definir até que ponto os dados de localização poderiam ser apresentados para ajudar na busca. Na pesquisa de campo, 16 dos 20 respondentes (80\%) preferiram que o endereço residencial e o número do imóvel não fossem exibidos ao público em geral (Apêndice B, Figura~B8). Esse retorno orientou a apresentação de uma localização aproximada nas telas públicas. Para facilitar o contato sem expor diretamente os dados pessoais, o aplicativo oferece mensagens privadas. O cartaz, por ser destinado ao compartilhamento externo, ainda pode conter dados fornecidos pelo autor da publicação, por isso sua divulgação exige atenção. A dificuldade foi conciliar a precisão necessária para localizar o animal com as preferências de privacidade apontadas pelos próprios participantes.

Outra dificuldade foi estabelecer uma pontuação útil para relacionar anúncios de animais perdidos e encontrados. O sistema soma ou subtrai pesos definidos manualmente para características como espécie, sexo, porte, raça, pelagem, idade e distância. Na versão atual, a combinação desses critérios tem produzido sugestões de correspondência com frequência maior do que o autor considera desejável, inclusive quando as semelhanças não são suficientes para uma identificação segura. Como os pesos e o limiar ainda não foram calibrados sistematicamente com exemplos de correspondências confirmadas e descartadas, a pontuação expressa apenas uma similaridade heurística, não a probabilidade de os anúncios tratarem do mesmo animal. Cabe aos usuários verificar as informações e confirmar ou descartar cada sugestão.

A integração entre esses recursos também exigiu atenção às permissões: dados de perfil, publicações e mensagens não têm todos o mesmo nível de acesso. Foram implementadas políticas de segurança em nível de linha (RLS), descritas no Capítulo 6, e sua configuração foi inspecionada manualmente no Capítulo 8. Essa inspeção, contudo, não equivale a uma auditoria independente nem a testes de invasão.

\subsection{Trabalhos Futuros}

Uma prioridade para a evolução do WeFIND é reduzir sugestões inadequadas na correspondência entre publicações. Para isso, os pesos e o limiar podem ser ajustados com um conjunto de pares de anúncios revisados por pessoas, identificando quais representam o mesmo animal e quais não correspondem. A partir desses exemplos, seria possível medir os acertos e os falsos positivos antes de alterar os critérios. Em seguida, pode-se avaliar uma funcionalidade de inteligência artificial que compare diretamente as fotografias das duas publicações, procurando semelhanças visuais que os campos preenchidos pelo usuário não descrevem. Essa análise deve complementar as informações de espécie, características e localização, sem confirmar automaticamente uma identidade. Sua contribuição precisaria ser testada em um conjunto de imagens representativo para verificar se realmente melhora a precisão das sugestões.

Também devem ser consideradas as questões de privacidade ao ampliar os recursos de localização e contato. Como próximos passos, podem ser avaliados controles mais claros sobre quais informações aparecem nas telas públicas, no código QR e nos cartazes, além de testes das permissões que protegem perfis, publicações e conversas.

Outra evolução possível é automatizar no backend a limpeza global de publicações expiradas. Na versão atual, a rotina é acionada localmente e considera apenas as publicações do usuário autenticado; por isso, a exclusão definitiva das publicações de outras contas depende de a rotina ser executada no contexto de cada usuário.

As respostas abertas dos participantes apontaram melhorias específicas de usabilidade: destacar melhor o acesso aos detalhes de uma publicação, explicar a diferença entre os tipos de registro, orientar quem usa o aplicativo pela primeira vez e facilitar o acesso ao histórico e às publicações salvas. Essas mudanças podem ser implementadas e avaliadas por meio de novas sessões em que a conclusão, as dúvidas e o tempo de cada tarefa sejam observados separadamente.

Além disso, uma nova avaliação com participantes de diferentes perfis poderia complementar o questionário SUS com observação direta das tarefas. Testes em aparelhos e condições de rede distintos, assim como uma avaliação dinâmica de segurança, também examinariam aspectos que não foram medidos neste trabalho.

Por fim, para investigar o impacto social do WeFIND, um estudo piloto em uma região delimitada poderia acompanhar publicações reais, com autorização dos envolvidos e critérios claros de registro. Indicadores como atualizações recebidas, tempo até o encerramento da publicação e reencontros informados ajudariam a descrever os resultados observados. Sem esse acompanhamento, a implementação das funcionalidades, por si só, não permite concluir que o aplicativo aumente a taxa de reencontros.
